\clearpage
\SourceLeaf{143}
\begin{center}
[III.]

{[ZUR GESCHICHTE DER ENTDECKUNG DER METHODE DER\\
KLEINSTEN QUADRATE.]}
\end{center}

\begin{center}
Allgemeine Geographische Ephemeriden. Herausgegeben von F.\ \textsc{von Zach}.\\
Vierter Band. 4. Stück. October 1799. S.\,378. Weimar, 1799.
\end{center}

Erlauben Sie, dass ich einen im Julius-St.\ der A.\,G.\,E. bemerkten Druckfehler anzeige. S.\,xxxv der Einleitung, bei der Angabe des Bogens zwischen dem Panthéon und Évaux, muss statt 76545,74 stehen 76145,74. Die Summe ist richtig, und der Fehler kann auf kein anderes als dieses Stück fallen\HistoricalNote{*)}{*)}{Dieser Druckfehler hat seine Richtigkeit, und ist auch aus dem beigesetzten Decimal-Grade $2^{0}\!,66868$ zu erkennen. v.\,Z.}. Ich entdeckte diesen Fehler, indem ich meine Methode, von der ich Ihnen eine Probe gegeben habe\HistoricalNote{**)}{**)}{Hiervon ein andermal. v.\,Z.}, anwandte, um bloss aus diesen vier gemessenen Stücken die Ellipse zu bestimmen, und $\frac{1}{150}$ Abplattung fand; nach Verbesserung jenes Fehlers fand ich $\frac{1}{187}$, und den ganzen Quadranten 2565006 Modulen (nemlich ohne Rücksicht auf den Grad in Peru). Der Unterschied von $\frac{1}{150}$ und $\frac{1}{187}$ ist in diesem Falle eben nicht erheblich, da die End-Punkte zu nahe liegen.

\bigskip
\noindent Braunschweig, den 24. Aug.\ 1799. \hfill C.\,F.\,\textsc{Gauss}.

